\chapter{La conformité multi-états}
\label{chap:multietats}

\begin{cle}
\textbf{Idée directrice.} Le besoin de capital dépend de l'état de conformité de l'entité,
objet ni binaire ni statique. On le modélise à deux échelles de temps (\emph{fast-slow}). À
l'échelle rapide, un Vasicek multi-états définit l'état (non conforme, partiellement conforme,
conforme) de chaque pilier et le fait piloter la cascade. À l'échelle lente, une chaîne de Markov
fait progresser la conformité le long de la mise en conformité. Tout ce qui en sort se lit avec
la règle que les chapitres~\ref{ch:identifiabilite} et~\ref{chap:partielle} établissent : le
niveau est illustratif, l'écart et l'ordre sont ce qui est défendable.
\end{cle}

\section{Le principe des deux échelles de temps}

Le moteur du chapitre~\ref{chap:cascade} traite l'état de conformité
comme \emph{figé}, et les niveaux qui en découlent ne sont établis qu'en
partie~\ref{part:resultats}. Or la
conformité n'est ni un point ni un instantané : elle est graduelle (trois niveaux au moins) et
évolue dans le temps. On distingue donc deux dynamiques. La couche rapide est celle des
incidents et de leur cascade à l'intérieur d'une année, chiffrée par le moteur du
chapitre précédent, \emph{conditionnellement} à une configuration de conformité. La couche
lente est celle de la gouvernance : l'état de conformité de chaque pilier progresse sur des
mois et des années. L'état de la couche lente sert de paramètre à la couche rapide.
C'est exactement la structure d'un modèle de risque de crédit à migration de rating, sur
laquelle on revient en \S\ref{sec:couplage}.

\paragraph{Cette séparation n'est pas une commodité de calcul, et un précédent le confirme.} On
pourrait la lire comme un artifice permettant de simuler deux dynamiques sans les coupler. Elle
répond en réalité à une exigence de fond : la propagation à l'intérieur d'un sinistre et
l'évolution de la conformité d'une année sur l'autre ne vivent pas sur la même horloge, et les
confondre reviendrait à laisser une tendance longue expliquer un enchaînement de quelques
heures, ou l'inverse. Le préprint de cascade climatique cité au chapitre~\ref{ch:positionnement}
\citep{ccrn2026} sépare de la même façon trois indices, un temps calendaire, une position
spatiale et un pas de propagation \emph{à l'intérieur} d'un événement, en donnant explicitement
pour motif d'empêcher que l'évolution de long terme ne soit confondue avec la propagation de
court terme. Une équipe travaillant sur un autre péril aboutit donc à la même séparation pour la
même raison, ce qui la déplace d'un choix de modélisation vers une contrainte de construction.

\section{Le Vasicek multi-états : latente et seuils ordonnés}

Chaque pilier porte une capacité de conformité latente $C^\star\sim\mathcal N(0,1)$,
décomposée à la Vasicek en un facteur systémique commun $\Theta$ et un choc idiosyncratique :
\begin{equation}
C^\star \;=\; \gamma\,\Theta \;+\; \sqrt{1-\gamma^2}\,\varepsilon, \qquad \gamma=0{,}68.
\end{equation}
Au lieu d'un unique seuil de défaut, on en pose deux, ordonnés, qui découpent trois
états :
\begin{equation}
\text{NC si } C^\star\le K_{\text{bas}}, \quad
\text{PC si } K_{\text{bas}}<C^\star\le K_{\text{haut}}, \quad
\text{C si } C^\star>K_{\text{haut}}.
\end{equation}
Les seuils sont fixés par les probabilités d'état ancrées sur les enquêtes de
conformité (fin 2025), $K_{\text{bas}}=\Phi^{-1}(p_{\text{NC}})$ et
$K_{\text{haut}}=\Phi^{-1}(p_{\text{NC}}+p_{\text{PC}})$. Ce dispositif est un
\emph{ordered-probit} posé sur la structure de Merton-Vasicek : il réutilise la machinerie de
seuil du chapitre~\ref{chap:cascade} (le seuil $K_j$) et généralise la variable latente de
défaut, binaire, en une variable latente de conformité à trois niveaux.

\begin{cle}
\textbf{Le facteur systémique rend les probabilités conditionnelles.} La probabilité d'être dans
chaque état sachant l'environnement $\Theta$ s'écrit
\begin{equation}
\mathbb P(\text{NC}\mid\Theta)=\Phi\!\Big(\tfrac{K_{\text{bas}}-\gamma\Theta}{\sqrt{1-\gamma^2}}\Big),
\qquad
\mathbb P(\text{C}\mid\Theta)=1-\Phi\!\Big(\tfrac{K_{\text{haut}}-\gamma\Theta}{\sqrt{1-\gamma^2}}\Big).
\end{equation}
Intégrée sur $\Theta$, la marginale redonne exactement l'ancrage (contrôle de cohérence) ; en
crise systémique, la masse bascule vers la non-conformité. C'est la version à trois états
ordonnés de l'amplification systémique d'un modèle à facteur unique.
\end{cle}

% SOURCES-SCRIPTS: 63


\section{Les quatre canaux : comment l'état pilote la cascade}
\label{sec:quatre-canaux}

L'état de conformité agit sur la cascade par quatre canaux, câblés sans réécrire
l'agrégation du chapitre~\ref{chap:cascade} :
\begin{enumerate}[leftmargin=1.5em,itemsep=2pt,topsep=2pt]
\item \textit{Fréquence} $\lambda$ : l'état module le nombre d'incidents amorcés, de $21{,}56$
  incidents matériels par an à l'état conforme à $53{,}62$ à l'état non conforme, par le
  multiplicateur de scénario décrit ci-dessous.
\item \textit{Propagation} $g$ : une entité plus conforme contient mieux la contagion, donc un
  gain plus faible, de $0{,}90$ à $0{,}45$. La propension à propager depuis un pilier reste
  $e_j=g(c_j)\,s_j/\max_k s_k$, avec un gain indexé sur l'état du pilier source.
\item \textit{Détection} $p_u$ : le taux de matérialité (le seuil de détection du
  chapitre~\ref{chap:cascade}). Une entité conforme intercepte et remédie tôt, de sorte qu'une
  moindre fraction d'incidents atteint la sévérité matérielle : $p_u$ est multiplié par $0{,}85$
  à l'état conforme et par $1{,}20$ à l'état non conforme, deux valeurs posées.
\item \textit{Accumulation par un prestataire commun} : à l'état non conforme, un incident amorcé
  sur le pilier des tiers devient un choc commun qui touche chacun des autres piliers avec la
  probabilité $0{,}68$, au lieu de se propager comme un nœud de cascade (\S\ref{sec:p4-accumulation}).
\end{enumerate}
On présente ces canaux en lectures emboîtées (fréquence seule ; $+$ propagation ; $+$
détection ; $+$ accumulation), ce qui rend visible la contribution marginale de chacun et
empêche qu'un canal domine en silence. Le chiffre de tête du chapitre~\ref{chap:resultats}
relâche les quatre.

\paragraph{D'où vient le multiplicateur de fréquence.} Il se compose vecteur d'attaque par
vecteur d'attaque : $\lambda_{\text{NC}}=\lambda_{\text{C}}\sum_v \pi_v\,m_v$, où $\pi_v$ est la
part du vecteur $v$ dans le recensement Hackmageddon (\S\ref{sec:hackmageddon}) et $m_v$ le centre
d'une fourchette de multiplicateur propre au scénario non conforme.

\begin{center}\footnotesize
\renewcommand{\arraystretch}{1.2}
\begin{tabular}{@{}l r c r r >{\raggedright\arraybackslash}p{5.2cm}@{}}
\toprule
\textbf{Vecteur} & \textbf{Part $\pi_v$} & \textbf{Fourchette} & \textbf{Centre} & \textbf{$\pi_v\,m_v$} & \textbf{Effet publié qui ancre la fourchette} \\
\midrule
Hameçonnage & $0{,}388$ & $[1{,}8\,;3{,}0]$ & $2{,}40$ & $0{,}9312$ & réduction du risque par la formation (Ponemon/IBM) \\
Exploitation de vulnérabilité & $0{,}338$ & $[2{,}0\,;2{,}6]$ & $2{,}30$ & $0{,}7774$ & écart de conversion en intrusion entre exploitation et hameçonnage (ENISA) \\
Chaîne d'approvisionnement & $0{,}158$ & $[1{,}9\,;2{,}6]$ & $2{,}25$ & $0{,}3555$ & surcoût et progression des brèches d'origine tierce (IBM, SecurityScorecard) \\
Identifiants compromis & $0{,}063$ & $[3{,}0\,;8{,}0]$ & $5{,}50$ & $0{,}3465$ & réduction de compromission par l'authentification forte (Microsoft Research) \\
Autres & $0{,}053$ & $[1{,}3\,;1{,}6]$ & $1{,}45$ & $0{,}0769$ & aucune source dédiée \\
\midrule
\multicolumn{4}{@{}l}{Multiplicateur agrégé, état non conforme} & $2{,}4875$ & $21{,}56\times 2{,}4875 = 53{,}62$ incidents par an \\
\bottomrule
\end{tabular}
\end{center}

\noindent Le multiplicateur est \emph{semi-ancré} et non calibré. Les parts sont une citation
externe, et les bornes des fourchettes traduisent chacune un effet publié en multiplicateur
prudent, sans estimation sur la donnée du projet : sa provenance est tracée, sa valeur est posée.
Deux conséquences sont à garder. Aux bornes des fourchettes, le multiplicateur va de $1{,}933$ à
$3{,}042$, soit une fréquence non conforme de $41{,}7$ à $65{,}6$ incidents par an. Et le vecteur
des identifiants, qui ne compte que $6{,}3\,\%$ des incidents, porte $13{,}9\,\%$ du multiplicateur,
sa fourchette étant la plus large.

% SOURCES-SCRIPTS: 67 99 43


\section{Les valeurs de \texorpdfstring{$g$}{g} sont posées, et la conclusion n'en dépend pas}
\label{sec:invariance-g}

Le gain de propagation $g$ est l'un des canaux par lesquels la conformité agit, et ses trois
valeurs,
\[
  g = 0{,}45 \ \text{(conforme)}, \qquad 0{,}68 \ \text{(partiellement conforme)},
  \qquad 0{,}90 \ \text{(non conforme)},
\]
sont posées. Elles ne sont calibrées sur aucune donnée, et pour une raison
irréductible : aucune entité ne publie son état de conformité \textsc{dora}, si bien qu'il
n'existe aucun échantillon d'où tirer la correspondance état $\to g$. Un écart énoncé comme
\og le capital passe de $116$ à $169$~M\euro\fg{} hérite donc entièrement de ces trois décimales.
Ce qui s'établit sans elles tient en deux temps.

\paragraph{Le lemme, qui n'allait pas de soi.} Le capital est croissant en $g$. Ce
n'est pas une évidence : le chapitre~\ref{chap:resultats} établit qu'à indice de queue élevé la
\textsc{var} annuelle est \emph{peu sensible} à la structure de cascade, par principe du grand
saut unique, et peu sensible ne veut pas dire monotone. On le vérifie donc numériquement plutôt
que de le supposer, sur une grille de vingt points et à nombres aléatoires communs, condition
sans laquelle deux valeurs voisines de $g$ ne différeraient que par le bruit de tirage. Résultat :
\textbf{aucun pas décroissant} sur les vingt intervalles, et cela sur quatre grandeurs à la fois,
le capital et la perte moyenne, à l'échelle d'entité comme à celle du secteur.

\paragraph{L'invariance, qui en découle.} Si le capital est croissant en $g$, alors toute
application état $\to g$ respectant l'ordre $g_C < g_{PC} < g_{NC}$ produit
$\mathrm{SCR}_C < \mathrm{SCR}_{PC} < \mathrm{SCR}_{NC}$, \emph{quelles que soient les valeurs}.
Or la charge forfaitaire de Formule Standard ne dépend pas de l'état. L'existence et le
\emph{sens} de l'écart sont donc invariants à la calibration ; seule son amplitude
est un scénario.

\begin{center}\small
\renewcommand{\arraystretch}{1.25}
\begin{tabular}{@{}l r r r r r@{}}
\toprule
\textbf{Application état $\to g$} & \textbf{Conf.} & \textbf{Part.} & \textbf{Non conf.} &
\textbf{Écart} & \textbf{Ordre} \\
 & M\euro & M\euro & M\euro & & \\
\midrule
étroite $(0{,}70 / 0{,}80 / 0{,}90)$ & $145{,}5$ & $156{,}5$ & $169{,}0$ & $+16{,}2\,\%$ & oui \\
\textbf{retenue} $(0{,}45 / 0{,}68 / 0{,}90)$ & $\mathbf{116{,}4}$ & $\mathbf{143{,}2}$ &
$\mathbf{169{,}0}$ & $\mathbf{+45{,}2\,\%}$ & oui \\
large $(0{,}20 / 0{,}55 / 0{,}90)$ & $94{,}4$ & $126{,}4$ & $169{,}0$ & $+79{,}1\,\%$ & oui \\
conformité faible $(0{,}45 / 0{,}55 / 0{,}65)$ & $116{,}4$ & $126{,}4$ & $139{,}2$ &
$+19{,}5\,\%$ & oui \\
très étroite $(0{,}85 / 0{,}875 / 0{,}90)$ & $162{,}6$ & $166{,}7$ & $169{,}0$ & $+3{,}9\,\%$ &
oui \\
\midrule
\emph{repère forfaitaire de Formule Standard} & $450$ & $450$ & $450$ & $0\,\%$ & \emph{plat} \\
\bottomrule
\end{tabular}
\end{center}

\vspace{2pt}
La dernière ligne du tableau est la plus instructive. Même en n'accordant à la conformité qu'un
déplacement de $0{,}05$ sur $g$, soit une hypothèse presque nulle, l'écart de capital subsiste
($+3{,}9\,\%$) quand le forfait reste exactement plat. La thèse du mémoire est cet écart de signe, non une
valeur de $g$.

\begin{cle}
\textbf{Ce qui reste une hypothèse, et c'est tout ce qui reste.} La proposition
\og se conformer à \textsc{dora} réduit la propagation entre piliers\fg{}, soit $g_C < g_{NC}$. C'est
une affirmation qualitative, discutable, et sur laquelle le règlement lui-même prend parti
puisque ses cinq piliers sont des exigences de maîtrise. Elle est d'un autre ordre de solidité
qu'un triplet de valeurs. Ce qui reste un scénario : l'amplitude, qui ne dépend que de
l'écartement $g_{NC}-g_C$. L'élasticité d'arc du capital au gain vaut $0{,}41$ sur
$g \in [0{,}20\,;0{,}95]$, de sorte qu'un lecteur qui préfère ses propres valeurs recalcule son
écart sans relancer le modèle. Et pour que cet écart dépasse $10\,\%$, il suffit que la
conformité fasse baisser $g$ de $0{,}90$ à $0{,}75$, soit un écartement de $0{,}15$ : l'exigence
porte sur un ordre de grandeur, pas sur une mesure.
\end{cle}

Une conséquence de rédaction en découle, et elle est appliquée au
chapitre~\ref{chap:resultats} : l'écart entre états n'y est pas présenté comme une mesure mais
comme un \emph{scénario d'écartement}, et c'est l'invariance d'ordre qui est mise en avant,
puisque c'est elle qui est établie (figure~\ref{fig:invariance-g}).

\begin{figure}[htbp]
\centering
\figover{S23_invariance_conformite.png}
\caption{La thèse ne dépend pas des valeurs de $g$, seulement de
leur ordre. (a)~le lemme de monotonie sur vingt points de grille, à nombres aléatoires
communs, avec les trois valeurs retenues situées sur la courbe et la charge forfaitaire qui,
elle, ne dépend pas de l'état. (b)~cinq applications état $\to g$, de très étroite à
large : l'ordre du capital est respecté par toutes, le forfait reste plat dans tous les cas.
(c)~l'amplitude de l'écart en fonction du seul écartement $g_{NC}-g_C$, la valeur
retenue étant marquée : c'est la seule grandeur qui soit un scénario.}
\label{fig:invariance-g}
\end{figure}
% SOURCES-SCRIPTS: 66 67

\section{L'état par pilier}

L'apport central est de porter l'état par pilier : chaque pilier a sa propre latente
$C^\star_j$, les cinq étant corrélées par le facteur systémique commun $\Theta$. C'est ce qui
fait \emph{interagir} la conformité et la cascade : un pilier non conforme propage plus et
détecte moins que ses voisins conformes. Les trois canaux deviennent indexés par pilier
($\lambda_j$, $g_j$, $p_{u,j}$) et alimentent un moteur d'agrégation qui se réduit exactement à
la version globale lorsque tous les piliers partagent un état (contrôle de cohérence). On en
déduit la décomposition du surcoût de non-conformité par pilier, en basculant un pilier
à la fois, dont les résultats sont au chapitre~\ref{chap:resultats}.

\subsection{La dépendance entre états : structure portée et structure omise}
\label{sec:dependance-etats}

% SOURCES-SCRIPTS: 69 77 94

Une objection métier se pose ici, et elle est juste : le
recouvrement des dispositifs de contrôle fait qu'être conforme aux piliers~1 et~2 confère
indirectement une conformité partielle au pilier~3, les mêmes preuves d'audit servant aux trois.
Le modèle ne représente pas ce mécanisme ; ce qu'il porte déjà et ce que l'omission coûte se
mesurent.

\textbf{Ce qu'il porte déjà.} Les états ne sont pas indépendants : les cinq latentes partagent le
facteur systémique avec une charge de $0{,}68$, soit une corrélation de $0{,}462$ entre deux
piliers quelconques. Répondre à l'objection par \og{}la dépendance manque\fg{} serait donc faux.
Ce qui manque est sa structure : la dépendance est \emph{échangeable}, identique pour
toutes les paires, alors que le recouvrement des contrôles est propre à certains triplets.

\textbf{Ce que l'omission coûte, et la réponse est nulle.} On ajoute un surcroît de corrélation
$\delta$ aux deux seules paires concernées, par une matrice de corrélation dont la diagonale
vaut un, ce qui préserve les marges \emph{par construction} et isole donc l'effet de structure.
Le surcroît admissible va jusqu'à $\delta = 0{,}38$ avant que la matrice cesse d'être définie
positive. Sur ce balayage, la loi des configurations bouge nettement : la probabilité que les
cinq piliers soient non conformes passe de $0{,}068$ à $0{,}088$ et celle qu'ils soient tous
conformes de $0{,}291$ à $0{,}330$. Le capital espéré, lui, ne bouge que de $2$~M\euro, soit
$0{,}02\,\%$.

\textbf{Et la raison n'est pas numérique, elle est structurelle.} Le capital des $32$
configurations s'ajuste presque exactement par une forme \emph{additive} en indicateurs de
pilier, $R^2 = 0{,}9945$, avec des contributions de $3\,524$ (P1), $2\,863$ (P4), $2\,022$ (P2),
$1\,577$ (P3) et $897$~M\euro{} (P5) au-dessus d'un socle tous conformes. Or l'espérance d'une
fonction additive ne dépend que des marges, jamais de la dépendance : les termes croisés
sont le seul canal par lequel une structure pourrait agir, et il n'en reste presque rien.
L'invariance ne tient donc pas à la valeur de $\delta$ retenue, elle vaut pour \emph{toute}
structure de dépendance à marges fixées, y compris celles qu'on n'a pas essayées. C'est ce qui
rend la conclusion solide alors qu'aucune donnée du projet ne calibre $\delta$. Le détail et la
figure sont déportés \matcomp.

Cette quasi-additivité est par ailleurs le même fait que l'interaction entre piliers déjà
publiée, petite devant la somme des marginaux et de signe non résolu (\S\ref{sec:allocation}) :
deux lectures d'un seul constat, et l'ordre des contributions additives reproduit celui des
marginaux, $P1 > P4 > P2 > P3 > P5$.

\paragraph{L'argument porte sur une espérance, et le capital est un quantile.} Il ne
transporte ni au quantile de la loi de perte mélangée sur les configurations, ni au quantile de
la loi des configurations elle-même. Sur le premier, l'invariance est seulement non détectée ;
sur le second, l'effet est de quelques pour cent, porté par une polarisation des configurations
extrêmes. Les deux calculs et leurs limites sont déportés \matcomp.

\section{La chaîne de Markov et les séjours de type phase}

À l'échelle lente, l'état de conformité suit une chaîne de Markov à temps continu
$\text{NC}\to\text{PC}\to\text{C}$, l'état Conforme étant absorbant en cas de base
(une fois conforme, on le reste ; la régression est un axe de sensibilité). Les séjours ne sont
pas exponentiels mais de type phase (Erlang-$k$) : une exponentielle autoriserait une
remédiation quasi instantanée (densité maximale en zéro), irréaliste, alors qu'un projet DORA a
une durée minimale. L'Erlang-$k$, somme de $k$ exponentielles, a un mode strictement positif qui
capte cette durée.

\begin{attention}
\textbf{Deux Erlang, deux objets sans rapport.} Le mot désigne ici la loi des \emph{durées de
séjour} dans la chaîne de conformité, où $k$ est un choix de modélisation assumé. Il désigne
ailleurs, au chapitre~\ref{chap:donnees}, la loi \emph{mélangeante} de la fréquence, dont la
forme ajustée vaut $0{,}63$ et n'est donc pas une Erlang, celle-ci exigeant une forme
entière supérieure ou égale à un. Les deux usages n'ont ni le même support, ni le même statut :
le premier est posé, le second est estimé et rejeté comme tel.
\end{attention} On développe chaque macro-état transitoire en $k$ phases et l'on résout le
générateur $Q$ sur l'espace de phases par exponentielle de matrice :
\begin{equation}
\big(\mathbb P(\text{NC},t),\mathbb P(\text{PC},t),\mathbb P(\text{C},t)\big)
\;=\; \text{(sommes de phases de)}\;\; v_0\,\exp(Qt),
\end{equation}
où $v_0$ est la distribution initiale. À chaque horizon $t$, cette distribution pondère les
besoins de capital par état du Vasicek multi-états :
\begin{equation}
\overline{\mathrm{SCR}}(t)\;=\;\sum_{e\in\{\text{NC,PC,C}\}}\mathbb P(e,t)\,\mathrm{SCR}_e .
\end{equation}
Cette grandeur est une \emph{moyenne de quantiles}, pondérée par la loi de l'état à l'horizon
$t$. Ce n'est pas le quantile à $99{,}5\,\%$ de la perte d'une entité dont l'état est incertain,
qui serait le quantile de la loi mélangée : la queue d'un mélange se prélève dans ses composantes
les plus lourdes, et les deux objets ne coïncident pas. Sur le mélange des $32$ configurations de
piliers, l'écart mesuré est de $9{,}3\,\%$ \matcomp. La
trajectoire se lit donc comme l'espérance du besoin de capital sous l'incertitude d'état, qui est
la convention retenue dans tout le mémoire.

% SOURCES-SCRIPTS: 08b 77


\section{Le couplage systémique et l'analogie de crédit}
\label{sec:couplage}

\begin{cle}
\textbf{Un seul facteur couple les deux couches.} Le même facteur systémique $\Theta$ agit à la
fois sur la couche lente et sur la couche rapide. Un environnement dégradé ($\Theta<0$) fait deux
choses avec le \emph{même} tirage : il ralentit la remédiation, via des taux de
transition $\mu(\Theta)=\mu_0\,e^{\beta_{\text{rem}}\Theta}$, et il durcit la cascade,
via $\mathrm{SCR}_e(\Theta)=\mathrm{SCR}_e\,e^{-\beta_{\text{casc}}\Theta}$. Les deux poussent le
capital dans le même sens : la corrélation entre conformité et sinistralité, qu'auraient perdue
des couches indépendantes, est restaurée. La dispersion de $\Theta$ engendre la bande
d'incertitude de la trajectoire.
\end{cle}

\noindent Cette architecture est celle, éprouvée, du risque de crédit à migration de
rating : la chaîne de Markov joue le rôle de la migration d'un \og rating de conformité DORA\fg{}
(la qualité change lentement, d'année en année), et le Vasicek multi-états celui du modèle de
défaut conditionnel au rating (rapide, dans l'année, pour un rating donné). Transposer
cette structure à deux étages au risque cyber n'est pas une invention fragile mais l'application
d'un cadre reconnu, ce qui en fait un argument de robustesse et non un pari.

Ce chapitre en donne la méthodologie complète ; les résultats chiffrés et les figures qui en
découlent sont rassemblés au chapitre~\ref{chap:resultats}.

% =====================================================================

% SOURCES-SCRIPTS: 63

